% This is a modified version of Springer's LNCS template suitable for anonymized MICCAI 2025 main conference submissions. 
% Original file: samplepaper.tex, a sample chapter demonstrating the LLNCS macro package for Springer Computer Science proceedings; Version 2.21 of 2022/01/12

\documentclass[runningheads]{llncs}
%
\usepackage[T1]{fontenc}
% T1 fonts will be used to generate the final print and online PDFs,
% so please use T1 fonts in your manuscript whenever possible.
% Other font encodings may result in incorrect characters.
%
\usepackage{amsmath}
\usepackage{amssymb}
\usepackage{multirow, booktabs}
\usepackage{graphicx,verbatim}
\usepackage{url}
% Used for displaying a sample figure. If possible, figure files should
% be included in EPS format.
%
% If you use the hyperref package, please uncomment the following two lines
% to display URLs in blue roman font according to Springer's eBook style:
%\usepackage{color}
%\renewcommand\UrlFont{\color{blue}\rmfamily}
%\urlstyle{rm}
%
\begin{document}
%
\title{Scaling In-Context Segmentation with Hierarchical Supervision}
\titlerunning{Scaling In-Context Segmentation with Hierarchical Supervision}
% If the paper title is too long for the running head, you can set
% an abbreviated paper title here
%
\begin{comment}  %% Removed for anonymized MICCAI submission
\author{First Author\inst{1}\orcidID{0000-1111-2222-3333} \and
Second Author\inst{2,3}\orcidID{1111-2222-3333-4444} \and
Third Author\inst{3}\orcidID{2222--3333-4444-5555}}
%
\authorrunning{F. Author et al.}
% First names are abbreviated in the running head.
% If there are more than two authors, 'et al.' is used.
%
\institute{Princeton University, Princeton NJ 08544, USA \and
Springer Heidelberg, Tiergartenstr. 17, 69121 Heidelberg, Germany
\email{lncs@springer.com}\\
\url{http://www.springer.com/gp/computer-science/lncs} \and
ABC Institute, Rupert-Karls-University Heidelberg, Heidelberg, Germany\\
\email{\{abc,lncs\}@uni-heidelberg.de}}

\end{comment}

\author{Anonymized Authors}  %% Added for anonymized MICCAI submission
\authorrunning{Anonymized Author et al.}
\institute{Anonymized Affiliations \\
    \email{email@anonymized.com}}
  
\maketitle              % typeset the header of the contribution
%
\begin{abstract}
In-context learning (ICL) enables medical image segmentation models to adapt to new anatomical structures from limited examples, reducing the clinical annotation burden. However, standard ICL methods typically rely on dense, global cross-attention, which scales poorly with image resolution. While recent approaches have introduced localized attention mechanisms, they often lack explicit supervision on the selection process, leading to redundant computation in non-informative regions. We propose a hierarchical framework that combines explicit patch selection with multi-level supervision. By supervising the model at both the patch and global levels, our approach learns to actively identify and attend only to the most informative anatomical regions. Evaluated on TotalSegmentator, our method achieves competitive in-domain CT segmentation accuracy while significantly reducing computational overhead for high-resolution volumes compared to full-context baselines. Training and evaluation code are available at \url{https://anonymous.4open.science/r/ic_segmentation}.

\keywords{In-Context Learning  \and Medical Image Segmentation}
% Authors must provide keywords and are not allowed to remove this Keyword section.

\end{abstract}

\section{Related Work}

\textbf{Global In-Context Learning.}

Early adaptations of ICL for segmentation, such as UniverSeg \cite{butoiUniverSegUniversalMedical2023}, utilize encoder-decoder architectures with dense cross-attention (CrossBlock) to fuse target and context features. Similarly, Iris \cite{gaoShowSegmentUniversal2025} encodes context pairs into global task embeddings.

\textit{Limitation:} These methods require processing the full resolution of both target and context images to compute attention maps, resulting in quadratic complexity $\mathcal{O}((HW)^2)$ that prohibits scaling to large inputs.

\textbf{Selective Attention and Retrieval.}

To address scalability, recent works have moved toward selective processing. Tyche \cite{rakicTycheStochasticInContext2025} introduces stochasticity to explore multiple potential masks, effectively attending to different "modes" of the output distribution. Other approaches, such as EICSeg \cite{xieEICSegUniversalMedical2026}, employ retrieval mechanisms to align features before processing.

\textit{Limitation:} While these methods focus attention on specific zones or features, this selection is often implicit : learned solely through the final segmentation loss without direct supervision on what to select. Consequently, they often still require encoding the entire image volume before the attention mechanism can discard irrelevant regions.

\textbf{Patch-based and Efficient Approaches.}

Direct patch-based methods like Medverse \cite{huMedverseUniversalModel2025a} use sliding window inference to handle high resolutions, while Patchwork \cite{reisertDeepNeuralPatchworks2022a} employs a hierarchical stacking strategy to iteratively sample uncertain regions.

\textit{Limitation:} Pure sliding windows treat patches independently, severing global context.


\textbf{Our Approach:} Unlike implicit attention methods, we employ explicit patch selection trained via multi-level supervision. This allows our model to discard uninformative background regions before heavy processing, retaining global context efficiency without the computational cost of full-volume attention.

\section{Method}

\noindent \textbf{Input:} Target image $I^t \in \mathbb{R}^{H \times W}$ and context set $\{(I^c_i, L^c_i)\}_{i=1}^{N_c}$ of image/label mask pairs.

\noindent \textbf{Output:} Predicted segmentation mask $\hat{L}^t$ at original resolution.

\subsection{Overview}

Our method processes images through a coarse-to-fine cascade of $M$ resolution levels $\{r_1, \ldots, r_M\}$ (e.g., $24 \rightarrow 96 \rightarrow 128$). At each level:
\begin{enumerate}
    \item \textbf{Sampling:} Select $K$ patches using Gumbel-top-$K$ sampling, guided by uncertainty from the previous level.
    \item \textbf{Attention:} Process target and context patches jointly through transformer layers.
    \item \textbf{Aggregation:} Combine patch predictions via averaging.
    \item \textbf{Refinement:} Additively fuse with upsampled previous-level predictions.
\end{enumerate}

\subsection{Differentiable Patch Sampling}

We use Gumbel-top-$K$ sampling to select patches in a differentiable manner:
\begin{equation}
    \text{idx} = \text{top-}K\left(\frac{w(x,y)}{\tau} + g_{x,y}\right), \quad g_{x,y} = -\log(-\log(u)), \; u \sim \text{Uniform}(0,1)
\end{equation}
where $w(x,y)$ are sampling weights and $\tau$ is a temperature parameter.

At level $\ell > 1$, sampling is guided by the entropy of the previous level's prediction:
\begin{equation}
    w_\ell(x, y) = H(p_{\ell-1}(x,y)), \quad H(p) = -p\log p - (1-p)\log(1-p)
\end{equation}
where $p_{\ell-1} = \sigma(\text{logits}_{\ell-1})$, focusing computation on uncertain regions.

Patches are sampled such that their \textbf{center} lies within the image, allowing partial overlap with boundaries. This increases coverage of border regions often undersampled by interior-only methods.

\subsection{Feature Extraction and Encoding}

Patches are encoded using frozen pretrained features from MedSAM2~\cite{maMedSAM2SegmentMedical2024}, a medical imaging foundation model. We extract multi-scale features from all three FPN levels (128$\times$128, 64$\times$64, 32$\times$32 at 512 input), resize to a common grid, and concatenate along the channel dimension, yielding 768-dimensional features per spatial location. A lightweight CNN encoder projects these to embeddings $\mathbf{z}_k \in \mathbb{R}^{d}$, producing skip connections for the decoder.

To enable a single backbone across all resolution levels, we condition on resolution via sinusoidal encoding, allowing generalization to unseen resolutions.

\subsection{Cross-Patch Attention}

Target and context patch embeddings are processed jointly through $L$ transformer layers with:
\begin{itemize}
    \item \textbf{Register tokens:} $R$ learnable tokens prepended to the sequence, providing global context aggregation across patches.
    \item \textbf{2D RoPE:} Rotary position embeddings encode spatial patch coordinates $(y, x)$, enabling reasoning about relative patch positions.
    \item \textbf{Type embeddings:} Per-layer learnable embeddings distinguish target from context patches.
    \item \textbf{Bidirectional attention:} All patches attend to all others (Flash Attention compatible).
\end{itemize}

A U-Net style CNN decoder reconstructs patch-level predictions using skip connections from the encoder.

\subsection{Aggregation and Level Combination}

Patch predictions are aggregated using simple averaging over overlapping regions:
\begin{equation}
    \hat{L}^t_\ell(x,y) = \frac{\sum_{k} \mathbf{1}_{(x,y) \in P_k} \cdot p_k(x,y)}{\sum_{k} \mathbf{1}_{(x,y) \in P_k}}
\end{equation}
where $\mathbf{1}_{(x,y) \in P_k}$ indicates pixel $(x,y)$ is covered by patch $k$.

Level predictions are combined via additive fusion within the coverage mask:
\begin{equation}
    \hat{L}^t_{\text{comb}} = \hat{L}^t_{\ell-1,\uparrow} + M_\ell \odot \hat{L}^t_\ell
\end{equation}
where $\hat{L}^t_{\ell-1,\uparrow}$ is the upsampled previous prediction and $M_\ell$ is the binary coverage mask indicating which pixels were sampled at level $\ell$.

\subsection{Training}

The model is trained with supervision on the final combined prediction:
\begin{equation}
    \mathcal{L} = \mathcal{L}_\text{combined}^M
\end{equation}
using soft Dice + Focal loss. During training, oracle sampling combines ground-truth foreground with entropy to balance coverage and uncertainty. \textbf{Oracle scheduling} linearly decays the oracle probability from 50\% to 0\% over the first epochs, transitioning to purely model-based (entropy-guided) sampling.

\begin{figure}[!htbp]
\centering
\includegraphics[width=\textwidth,keepaspectratio]{imgs/figures/architecture.pdf}
\caption{Architecture}
\label{fig:architecture}
\end{figure}

\begin{figure}[!htbp]
\centering
\includegraphics[width=\textwidth,keepaspectratio]{imgs/figures/attention.pdf}
\caption{Patch-level attention}
\label{fig:patch_attn}
\end{figure}

\section{Experiments}

\subsection{Dataset and Preprocessing}

We train and evaluate on TotalSegmentator~\cite{wasserthalTotalSegmentatorRobustSegmentation2023}, a large-scale CT dataset with 104 anatomical structures. To test true in-context generalization, we split by \textbf{label} rather than by case: 50 structures for training, 50 held-out structures for validation (Table~\ref{tab:label_splits}). This ensures the model must generalize to entirely unseen anatomical categories at test time, not just unseen patients.

\textbf{2D Slice Extraction.} We extract 2D slices from 3D volumes along all three anatomical axes (axial, coronal, sagittal), increasing data diversity and providing multi-view perspectives of each structure. Slices are filtered by coverage: only slices with $\geq 30$ foreground pixels are retained, preventing training on uninformative near-empty slices.

\textbf{Class-Balanced Sampling.} To prevent large organs (e.g., liver, lungs) from dominating training, we sample labels uniformly at each step. Within each label, we subsample to at most 6 slices per (case, axis) group using strided selection for even spatial distribution.

\textbf{Context Selection.} Context examples are preferentially drawn from the \textbf{same patient} as the target, providing anatomically consistent support. When insufficient same-case examples exist, we fall back to other patients with the same structure. We use $N_c = 5$ context pairs during training.

\textbf{Evaluation.}
\begin{itemize}
    \item \textbf{In-domain:} TotalSegmentator CT (validation labels, all cases)
    \item \textbf{Out-of-domain:} TotalSegmentator MRI~\cite{dantonoliTotalSegmentatorMRIRobust2025} (all labels, all cases)
\end{itemize}




\section{Results}
\subsection{In-Domain CT Segmentation}

\begin{table}[t]
\centering
\caption{Dice scores on TotalSegmentator CT validation set with 3 context pairs (mean $\pm$ std). Best per category in \textbf{bold}.}\label{tab:dice_totalseg_ct}
\begin{tabular}{lcc}
\toprule
Category & PatchICL & UniverSeg \\
\midrule
Organs (Thorax/Head/Spine)    & \textbf{0.460 $\pm$ 0.340} & 0.436 $\pm$ 0.375 \\
Muscles                       & 0.405 $\pm$ 0.297 & \textbf{0.414 $\pm$ 0.333} \\
Organs (Abd./Pelvis)          & 0.402 $\pm$ 0.291 & \textbf{0.437 $\pm$ 0.336} \\
Other                         & 0.345 $\pm$ 0.272 & \textbf{0.368 $\pm$ 0.313} \\
Vessels                       & 0.261 $\pm$ 0.288 & \textbf{0.305 $\pm$ 0.321} \\
\midrule
\textbf{Overall}              & 0.349 $\pm$ 0.286 & \textbf{0.375 $\pm$ 0.324} \\
\bottomrule
\end{tabular}
\end{table}

\textbf{Per-class analysis.} PatchICL outperforms UniverSeg on 19 of 58 validation classes and wins the Organs (Thorax/Head/Spine) category (+2.4\%). Structure-wise, PatchICL shows strong gains on pelvic organs (prostate +14\%, kidney cysts +17\%), bones (+1.9\% avg, femur +9\%), and mid-thoracic vertebrae (T6--T9: +5\% avg). On medium-difficulty structures (UniverSeg Dice 0.3--0.5), PatchICL matches or exceeds UniverSeg (+0.7\% avg). The gap widens on thin tubular structures: iliac vessels ($-$8 to $-$12\%) and lower ribs ($-$8\%).

\subsection{Cross-Modality Generalization (MRI)}

To evaluate cross-modality generalization, we test on TotalSegmentator MRI~\cite{dantonoliTotalSegmentatorMRIRobust2025} without any fine-tuning. Both methods were trained exclusively on CT data.

\begin{table}[t]
\centering
\caption{Dice scores on TotalSegmentator MRI with 3 context pairs (mean $\pm$ std). Models trained on CT only. Best per category in \textbf{bold}.}\label{tab:dice_totalseg_mri}
\begin{tabular}{lcc}
\toprule
Category & PatchICL & UniverSeg \\
\midrule
Organs (Thorax/Head/Spine)    & 0.641 $\pm$ 0.308 & \textbf{0.697 $\pm$ 0.312} \\
Muscles                       & \textbf{0.488 $\pm$ 0.282} & 0.465 $\pm$ 0.301 \\
Other                         & 0.484 $\pm$ 0.283 & \textbf{0.515 $\pm$ 0.307} \\
Organs (Abd./Pelvis)          & 0.415 $\pm$ 0.294 & \textbf{0.437 $\pm$ 0.321} \\
Bones (Spine)                 & 0.400 $\pm$ 0.210 & \textbf{0.449 $\pm$ 0.256} \\
Bones (Limbs/Pelvis)          & 0.321 $\pm$ 0.309 & \textbf{0.376 $\pm$ 0.348} \\
Vessels                       & 0.304 $\pm$ 0.271 & \textbf{0.379 $\pm$ 0.318} \\
\midrule
\textbf{Overall}              & 0.424 $\pm$ 0.286 & \textbf{0.455 $\pm$ 0.313} \\
\bottomrule
\end{tabular}
\end{table}

\textbf{Cross-modality analysis.} Despite training exclusively on CT, PatchICL wins the Muscles category (+2.3\%), with gains on autochthon (+4.5\%), gluteus maximus (+3.4\%), and stomach (+3.6\%). The gap narrows compared to in-domain CT ($-$3.1\% vs $-$2.6\% overall).

\subsection{Out-of-Domain Generalization (MedSegBench)}

To evaluate generalization beyond CT imaging, we test on MedSegBench~\cite{kusMedSegBenchComprehensiveBenchmark2024}, a comprehensive benchmark spanning 35 datasets across diverse modalities including ultrasound, X-ray, endoscopy, dermoscopy, fundus imaging, and microscopy. Both methods were trained exclusively on TotalSegmentator CT.

\begin{table}[t]
\centering
\caption{Dice scores on MedSegBench with 3 context pairs (mean $\pm$ std). Models trained on CT only.}\label{tab:dice_medsegbench}
\begin{tabular}{lccc}
\toprule
Difficulty & N & PatchICL & UniverSeg \\
\midrule
Hard (Dice $<$ 0.3)      & 47 & \textbf{0.110} & 0.088 \\
Medium (0.3--0.5)        & 12 & 0.286 & \textbf{0.383} \\
Good (0.5--0.7)          & 11 & 0.437 & \textbf{0.606} \\
Easy ($>$ 0.7)           & 15 & 0.601 & \textbf{0.800} \\
\midrule
\textbf{Overall (85 classes)} & -- & 0.412 $\pm$ 0.222 & \textbf{0.622 $\pm$ 0.287} \\
\bottomrule
\end{tabular}
\end{table}

\textbf{Analysis by difficulty.} Table~\ref{tab:dice_medsegbench} stratifies results by task difficulty (measured by UniverSeg baseline performance). PatchICL shows an advantage on the hardest tasks ($+$2.2\%), where both methods struggle to exceed 10\% Dice. On easier tasks where in-context cues are sufficient, UniverSeg's dense cross-attention provides more reliable transfer.

\textbf{Analysis by modality.} PatchICL achieves strong gains on ultrasound datasets (AbdomenUS: $+$13--24\% on multiple organs) and endoscopy (Kvasir polyps: $+$9.8\%, M2caiSeg surgical instruments: $+$7.8\%). These modalities exhibit high spatial variability where selective patch attention may focus on diagnostically relevant regions. Conversely, UniverSeg dominates on standardized protocols: dental X-ray (PanDental: $-$24\%), blood cell microscopy (WBC: $-$29\%), and fundus imaging (IDRIB: $-$37\%), where consistent imaging conditions favor global feature matching.

\subsection{Generalization to unseen dimensions}
Accuracy and flops
\subsubsection{Input image size}

\begin{figure}[!htbp]
\centering
\includegraphics[width=\textwidth]{imgs/plots/pareto_input_size.pdf}
\caption{Accuracy and compute by input dimensions}
\label{fig:dice_flops_per_input_size}
\end{figure}

\subsubsection{Number of context pairs}

\begin{figure}[!htbp]
\centering
\includegraphics[width=\textwidth]{imgs/plots/pareto_context_size.pdf}
\caption{Accuracy and compute by number of context pairs}
\label{fig:dice_flops_per_nb_context}
\end{figure}


\subsection{Explainability}
\paragraph{Attention}

\begin{itemize}
    \item Do register tokens correlate with label id, label size, target-context distance ? 
    \item Do empty patches attend to zeroAttn ? 
\end{itemize}

\begin{figure}[!htbp]
\centering
\includegraphics[width=\textwidth]{imgs/plots/reg_tokens_category.png}
\caption{2D UMAP projection of register 1, labeled by class category (TotalSegmentator, validation set)}
\label{fig:reg_tokens_label}
\end{figure}

\begin{figure}[!htbp]
\centering
\includegraphics[width=\textwidth]{imgs/plots/reg_tokens_class.png}
\caption{2D UMAP projection of register 1, labeled by class for each category (TotalSegmentator, validation set)}
\label{fig:reg_tokens_category}
\end{figure}

\paragraph{Explicit Sampling and Aggregation strategies}

Since we operate on every level explicitly, we can explore multiple sampling and aggregation strategies.

\begin{itemize}
    \item best sampling strategy ? Sample at highest probs ? Or most uncertain is better ? 
    \item best aggregation strategy ? replace previous probs everywhere ? Or only when certain ? 
\end{itemize}
 

\bibliographystyle{splncs04}
\bibliography{ref}

\newpage

\begin{table}[t]
\centering
\caption{Train/validation label split across 100 TotalSegmentator classes. Labels are split 50/50 while maintaining anatomical diversity in both sets.}\label{tab:label_splits}
\small
\begin{tabular}{|p{3.2cm}|p{5.5cm}|p{5.5cm}|}
\toprule
\textbf{Category} & \textbf{Training (50)} & \textbf{Validation (50)} \\
\midrule
Organs (Abd./Pelvis) & Duodenum, esophagus, R.~kidney & Colon, gallbladder, L.~kidney, liver, pancreas, prostate, small bowel, spleen, stomach, urinary bladder \\
\midrule
Organs (Thorax/ Head/Spine) & Brain, heart, L.~lung (upper), R.~lung (middle) & L.~lung (lower), R.~lung (lower), R.~lung (upper), spinal cord, thyroid gland, trachea \\
\midrule
Bones (Spine) & L1, L2, L4, L5, T2, T4, T7, T9, T10, T12 & C7, T1, T3, T5, T6, T8, T11, L3, S1, sacrum \\
\midrule
Bones (Ribs/ Sternum) & Sternum, L.~ribs 1--4, 7; R.~ribs 5--10 & Costal cart., L.~ribs 5, 6, 8--11; R.~ribs 1--4, 11 \\
\midrule
Bones (Limbs/ Shoulder/Pelvis) & Skull, L.~femur, L.~humerus, L.~scapula, R.~clavicula, R.~hip, R.~scapula & L.~clavicula, L.~hip, R.~femur, R.~humerus \\
\midrule
Muscles & L.~autochthon, R.~autochthon, L.~glut.~max., L.~glut.~med., L.~glut.~min., R.~glut.~min. & L.~iliopsoas, R.~iliopsoas, R.~glut.~max., R.~glut.~med. \\
\midrule
Vessels & IVC, SVC, L.~brachioceph.~v., L.~iliac~v., R.~iliac~v., L.~subclavian~a., R.~iliac~a., portal/splenic~v. & Aorta, L.~iliac~a., pulmonary~v., R.~subclavian~a. \\
\bottomrule
\end{tabular}
\end{table}

\newpage

\section{Important Messages to MICCAI Authors}

This is a modified version of Springer's LNCS template suitable for anonymized MICCAI 2026 main conference submissions. The author section has already been anonymized for you.  You cannot remove the author section to gain extra writing space.  You must not remove the abstract and keyword sections in your submission.

To avoid an accidental anonymity breach, you are not required to include the \textbf{Acknowledgments} and the \textbf{Disclosure of Interest} sections for the initial submission phase.  If your paper is accepted, you must reinsert these sections in your final paper.  If you opted to include these sections in the initial submission phase, they will count toward the 8-page limit.

To avoid desk rejection of your paper, you are encouraged to read "Avoiding Desk Rejection" by the MICCAI submission platform manager.  A few important rules are summarized below for your reference:  
\begin{enumerate}
    \item Do not modify or remove the provided anonymized author section except for inserting your paper ID.
    \item Do not remove the keyword section to gain extra writing space. Authors must provide at least one keyword.
    \item Do not remove the abstract sections to gain extra writing space. 
    \item Inline figures (wrapping text around a figure) are not allowed.
    \item Authors are not allowed to change the default margins, font size, font type, and document style.  If you are using any additional packages, it is the author's responsibility to ensure they do not inherently alter the document's layout.
    \item It is prohibited to use commands such as \verb|\vspace| and \verb|\hspace| to reduce the pre-set white space in the document.  Shrinking the white space between your figures/tables and their captions, and between sections and subsections, is strictly prohibited.
    \item Please make sure your figures and tables do not span into the margins.
    \item If you have tables, the smallest font size allowed is 8pt.  Your tables should not be inserted as figures.
    \item Please ensure your paper is properly anonymized.  Refer to your previously published paper in the third person.  Data collection location of private dataset should be masked.
        \item  If you are including a link to your code repository or datasets or submitting supplementary materials, you must ensure your repository and supplementary materials are anonymized.
    \item Your paper must not exceed the page limit.  You can have 8 pages of the main content (text (including paper title, author section, abstract, keyword section), figures and tables, conclusion, acknowledgment and disclaimer) plus up to 2 pages of references.  The reference section does not need to start on a new page.  
    \item *Beginning in 2025* 
    Supplementary materials are limited to multimedia content (e.g., videos) as warranted by the technical application (e.g., robotics, surgery, ...). These files should not display any proofs, analysis, additional results, or embedded slides, and should not show any identification markers either. Violation of this guideline will lead to desk rejection. PDF files may not be submitted as supplementary materials in 2026 unless authors are citing a paper that has not yet been published. In such a case, authors are required to submit an anonymized version of the cited paper. ACs will perform checks to ensure the integrity.
\end{enumerate}


\section{Section}
\subsection{A Subsection Sample}
Please note that the first paragraph of a section or subsection is
not indented. The first paragraph that follows a table, figure,
equation etc. does not need an indent, either.

Subsequent paragraphs, however, are indented.

\subsubsection{Sample Heading (Third Level)} Only two levels of
headings should be numbered. Lower level headings remain unnumbered;
they are formatted as run-in headings.

\paragraph{Sample Heading (Fourth Level)}
The contribution should contain no more than four levels of
headings. Table~\ref{tab1} gives a summary of all heading levels.

\begin{table}
\caption{Table captions should be placed above the
tables.}\label{tab1}
\begin{tabular}{|l|l|l|}
\hline
Heading level &  Example & Font size and style\\
\hline
Title (centered) &  {\Large\bfseries Lecture Notes} & 14 point, bold\\
1st-level heading &  {\large\bfseries 1 Introduction} & 12 point, bold\\
2nd-level heading & {\bfseries 2.1 Printing Area} & 10 point, bold\\
3rd-level heading & {\bfseries Run-in Heading in Bold.} Text follows & 10 point, bold\\
4th-level heading & {\itshape Lowest Level Heading.} Text follows & 10 point, italic\\
\hline
\end{tabular}
\end{table}


\begin{table}
\caption{Test table.}\label{tab2}
%\scalebox{0.8}
{
\begin{tabular}{|l|l|l|}
\hline
Heading level &  Example & Font size and style\\
\hline
\small small &  {\Large\bfseries Lecture Notes} & 14 point, bold\\
\tiny tiny &  {\large\bfseries 1 Introduction} & 12 point, bold\\
\normalfont normal & {\bfseries 2.1 Printing Area} & 10 point, bold\\
\fontsize{8}{9} 8pt & {\bfseries Run-in Heading in Bold.} Text follows & 10 point, bold\\
\scriptsize script & {\itshape Lowest Level Heading.} Text follows & 10 point, italic\\
\hline
\end{tabular}}
\end{table}

\noindent Displayed equations are centered and set on a separate
line.
\begin{equation}
x + y = z
\end{equation}
Please try to avoid rasterized images for line-art diagrams and
schemas. Whenever possible, use vector graphics instead (see
Fig.~\ref{fig1}).

\begin{figure}
\includegraphics[width=\textwidth]{fig1.eps}
\caption{A figure caption is always placed below the illustration.
Please note that short captions are centered, while long ones are
justified by the macro package automatically.} \label{fig1}
\end{figure}

\begin{theorem}
This is a sample theorem. The run-in heading is set in bold, while
the following text appears in italics. Definitions, lemmas,
propositions, and corollaries are styled the same way.
\end{theorem}
%
% the environments 'definition', 'lemma', 'proposition', 'corollary',
% 'remark', and 'example' are defined in the LLNCS documentclass as well.
%
\begin{proof}
Proofs, examples, and remarks have the initial word in italics,
while the following text appears in normal font.
\end{proof}
For citations of references, we prefer the use of square brackets
and consecutive numbers. Citations using labels or the author/year
convention are also acceptable. The following bibliography provides
a sample reference list with entries for journal
articles~\cite{ref_article1}, an LNCS chapter~\cite{ref_lncs1}, a
book~\cite{ref_book1}, proceedings without editors~\cite{ref_proc1},
and a homepage~\cite{ref_url1}. Multiple citations are grouped
\cite{ref_article1,ref_lncs1,ref_book1},
\cite{ref_article1,ref_book1,ref_proc1,ref_url1}.

 %% removed for anonymized MICCAI submission.
    
    % The following acknowledgement and disclaimer sections can be removed for the double-blind review process.  If and when your paper is accepted, reinsert the acknowledgement and the disclaimer clause in your final camera-ready version.
    % IF you opted to include the acknowledgement and disclaimer sections, they will count towards the 8-page limit.

\begin{credits}
\subsubsection{\ackname} A bold run-in heading in small font size at the end of the paper is
used for general acknowledgments, for example: This study was funded
by X (grant number Y).

\subsubsection{\discintname}
It is now necessary to declare any competing interests or to specifically
state that the authors have no competing interests. Please place the
statement with a bold run-in heading in small font size beneath the
(optional) acknowledgments\footnote{If EquinOCS, our proceedings submission
system, is used, then the disclaimer can be provided directly in the system.},
for example: The authors have no competing interests to declare that are
relevant to the content of this article. Or: Author A has received research
grants from Company W. Author B has received a speaker honorarium from
Company X and owns stock in Company Y. Author C is a member of committee Z.
\end{credits}


%
% ---- Bibliography ----
%
% BibTeX users should specify bibliography style 'splncs04'.
% References will then be sorted and formatted in the correct style.
%
% \bibliographystyle{splncs04}
% \bibliography{mybibliography}
%
\begin{thebibliography}{8}
\bibitem{ref_article1}
Author, F.: Article title. Journal \textbf{2}(5), 99--110 (2016)

\bibitem{ref_lncs1}
Author, F., Author, S.: Title of a proceedings paper. In: Editor,
F., Editor, S. (eds.) CONFERENCE 2016, LNCS, vol. 9999, pp. 1--13.
Springer, Heidelberg (2016). \doi{10.10007/1234567890}

\bibitem{ref_book1}
Author, F., Author, S., Author, T.: Book title. 2nd edn. Publisher,
Location (1999)

\bibitem{ref_proc1}
Author, A.-B.: Contribution title. In: 9th International Proceedings
on Proceedings, pp. 1--2. Publisher, Location (2010)

\bibitem{ref_url1}
LNCS Homepage, \url{http://www.springer.com/lncs}, last accessed 2023/10/25
\end{thebibliography}
\end{document}
